%---------------------------------------------------------
\section{Overview}
\label{sec:soa_overview}

The representation of a discretely sampled trajectory as an element of a
function space, obtained by expanding it in a finite system of basis
functions, is the operation that everything else in functional data
analysis is built on; \cite{ramsay2005} remains the reference work that
establishes this operation and the analyses constructed on top of it,
including FPCA. Chapter~\ref{chap:intro} argued that, in applications to
volatile economic series, this operation is rarely treated as an object
of empirical scrutiny in its own right, and Chapter~\ref{chap:theory}
developed the operator-theoretic machinery that any basis expansion feeds
into. This chapter reviews the literature that bears on that argument. A
systematic search of the bibliographic databases, screened against
explicit eligibility criteria, yielded 91 sources retained for full
analysis, of which 44 qualify as direct antecedents: work that both
justifies its choice of basis and evaluates the quality of the resulting
representation.

Section~\ref{sec:soa_basis_choice} reviews how functional basis families
are chosen and justified in practice. Section~\ref{sec:soa_fpca} examines
how that choice propagates into functional principal component analysis.
Section~\ref{sec:soa_smoothing} turns to smoothing and penalization, and
to the extent to which their effect on the fitted representation is
disentangled from the effect of the basis family itself.
Section~\ref{sec:soa_wavelets} reviews the use of wavelet bases for
non-smooth or non-stationary processes. Section~\ref{sec:soa_economics}
narrows the scope to economic and financial applications, with
particular attention to Latin American and Ecuadorian antecedents.
Section~\ref{sec:soa_gap} consolidates these five threads into a single
methodological gap, and Section~\ref{sec:soa_positioning} positions the
present thesis relative to it.

%---------------------------------------------------------
\section{Basis choice: how functional basis families are justified in
practice}
\label{sec:soa_basis_choice}

The reviewed literature converges markedly on B-spline bases as the
default choice for representing functional data: among the 91 reviewed
works, 30 of those that adopt a single basis family choose B-splines,
against 5 that choose Fourier and only 2 that choose wavelets. This
numerical convergence does not reflect a methodological consensus;
authors within the field state the underlying tension explicitly.
\cite{basna_driven_2022} describe the prevailing practice as ``a rather
ad hoc choice'' among Fourier, wavelet, or spline bases, and
\cite{eslami_spline_2024} report that ``a notable gap persists in the
literature regarding comprehensive comparisons between these
methodologies'' -- two works in this same corpus, two years apart, state
nearly the same objection: the choice of basis responds to convention
more often than to an explicit justification.

When a theoretical justification is offered, it typically rests on a
structural property attributed to the data, not on an empirical
comparison. \cite{shackleton_comparing_2024} are explicit on this point:
they prefer B-splines over Fourier ``due to the aperiodicity of the
data,'' while the general literature reserves Fourier bases precisely
for the opposite case, periodic signals \cite{nassar_empirically_2021}.
Wavelet bases, in turn, are justified not by periodicity but by
locality: \cite{yang_feature_2022} and \cite{amato_unsupervised_2025}
adopt them specifically to capture discontinuities or localized features
that a global basis over-smooths. This correspondence -- periodicity to
Fourier, global smoothness to B-splines, locality or discontinuity to
wavelets -- is the theoretical argument available in the field, but it
is rarely applied systematically: most of the 91 reviewed works simply
adopt B-splines without discussing whether the series at hand satisfies
the condition that would, in theory, favor that choice.

The minority of works that directly compare basis families (25 of 91)
shows that this choice is not cosmetic: results depend strongly on the
application domain and do not generalize from one setting to another.
\cite{kayano_principal_2009} find that a regularized Gaussian radial
basis outperforms B-splines in mean squared error for unbalanced data;
\cite{perezplaza_network_2018} find the opposite for GPS geodetic
curves, where P-splines outperform Fourier in two of three components;
\cite{amato_unsupervised_2025} report an even larger gap in spectral
classification, where a wavelet-based method reaches an adjusted Rand
index of 0.901 against only 0.320 for $k$-means on B-spline bases. This
heterogeneity of outcomes -- no single family dominates universally --
is precisely what makes it untenable to extrapolate the suitability of
B-splines, or of any other family, from one domain to another without
domain-specific evidence: if the choice of basis were immaterial, these
25 comparative works would not report different winners across
applications.

A second, more recent line of work (2016--2024) responds to the same
problem not by choosing more carefully among the classical families but
by avoiding the choice altogether: the matrix-completion methods of
\cite{descary_matrix_2016,descary_matrix_2019} and the orthonormal
``splinet'' bases of
\cite{nassar_empirically_2021,podgorski_splinets_2021,basna_spline_2024}
estimate an orthonormal basis directly from the data instead of fixing
one in advance. \cite{aguilera_comparative_2013}, in a register closer
to that of this thesis, show that penalizing the basis (P-splines)
matters as much as, or more than, the family chosen, reducing the mean
squared error substantially relative to unpenalized regression splines
within the same B-spline family -- a result that already anticipates
Section~\ref{sec:soa_smoothing}: the effect of the basis family and the
effect of the smoothing regime are not the same question, and the
general FDA literature is beginning to separate them. What it does not
do, as Section~\ref{sec:soa_economics} shows, is separate them for
economic or financial series.

%---------------------------------------------------------
\section{FPCA: how basis choice propagates into estimation and
truncation}
\label{sec:soa_fpca}

The most direct evidence that basis choice conditions the outcome of
FPCA does not come from a comparison of basis families in itself, but
from a controlled experiment that precedes it: \cite{tarpey_linear_2007}
fits the same curves with B-spline, Fourier, and power bases, fixing the
same dimension $p=5$ in all three cases to isolate the effect of the
family, and finds that ``$k$-means clustering results vary depending on
how the curves are fit to the data.'' This is the cleanest empirical
check in the corpus that PR2 has an affirmative answer: the basis family
not only affects the fit, it affects the downstream analysis built on
top of that fit. A second source of variation, orthogonal to the family,
lies in the formulation of FPCA itself: \cite{hormann_dynamic_2012} show
that a dynamic FPCA, computed in the frequency domain, explains more
variance than a static FPCA under the same underlying basis (80\% against
73\% in their application to PM10 pollution data) -- how the FPCA itself
is computed competes with the basis in which the curve is represented as
a source of differences in the results.

Where Section~\ref{sec:soa_basis_choice} found a partial convergence
toward B-splines in the choice of family, the criterion used to truncate
that basis -- how many basis functions to retain -- shows no comparable
convergence. The corpus documents at least eight distinct approaches in
active use: GCV, AIC/BIC and their variants, generalized information
criteria (GIC, MAIC, GBIC), automatic Bayesian selection (Bernoulli-type
latent variables in \cite{sousa_bayesian_2024}; sparse or variational
priors in \cite{tao_integrated_2025}), an eigenvalue-ratio criterion, as
used by \cite{linshang_nonstationary_2025}, cumulative explained-variance
thresholds, and purpose-built bootstrap procedures \cite{diks_noise_2023}.
\cite{beyaztas_comparison_2022} directly compare four of these criteria
(GCV, GIC, MAIC, GBIC) within the same study and treat the choice of
criterion as a design decision as consequential as the basis family
itself -- a reminder that choosing the basis well and choosing how many
functions of that basis to use are two independent decisions that are
rarely studied separately.

One antecedent isolates that independence explicitly, in a way that
directly anticipates the central argument of this thesis:
\cite{gao_factor_2024} deliberately fix a large dimension $K$ in their
B-spline basis and control smoothness exclusively through the
penalization parameter, instead of selecting $K$ by a criterion such as
GCV. It is, within the corpus, the methodological design closest to
separating the effect of basis dimension from the effect of the
smoothing regime -- but it does so within a single family (B-splines),
not as a comparison across families, and not on volatile economic series
(Section~\ref{sec:soa_economics}). This same tension, between selecting
the dimension automatically or decoupling it from smoothing, appears
again, in a different form, in Section~\ref{sec:soa_smoothing}.

%---------------------------------------------------------
\section{Smoothing, penalization, and parameter selection}
\label{sec:soa_smoothing}

The separation between the effect of the basis family and the effect of
the smoothing regime, which PR3 poses as an open question, does have
antecedents within a single family. \cite{aguilera_comparative_2013} hold
the B-spline family fixed and vary only the type of penalization -- none
(regression splines), continuous (smoothing splines), and discrete
(P-splines) -- finding that both penalized approaches reduce the mean
squared error substantially relative to the unpenalized case.
\cite{gao_factor_2024} take this separation further: they fix a large
dimension $K$ and control smoothness exclusively through the
penalization parameter, instead of selecting $K$ by a criterion such as
GCV. In both cases the basis family is held constant; what varies is the
smoothing regime, and the effect is measurable and substantial. This is
exactly the type of design that PR3 calls for, but applied within a
single family, not as a comparison across distinct families, and not on
economic or financial series.

When the practical selection of the smoothing parameter $\lambda$ is
examined, the most commonly cited automatic criterion, generalized
cross-validation, does not always perform as expected.
\cite{zin_smoothing_2020} document a case in which neither GCV nor
ordinary cross-validation produced a reasonable value of $\lambda$, and
the authors had to fall back on a subjective assessment of the fit. This
result is relevant to PR3 because it marks a practical limit: if the
automatic criterion for selecting $\lambda$ can fail, any comparison
between basis families that relies on an automatic, unverified selection
of the smoothing parameter risks confounding the effect of the family
with an artifact of a poor choice of $\lambda$.
\cite{beyaztas_comparison_2022} come closest, within the corpus, to
controlling both sources of variation at once: they compare three basis
families and, for each, four distinct selection criteria (GCV, GIC,
MAIC, GBIC), which allows the analysis to separate how much of the
observed difference comes from the family and how much from the
selection criterion. The application domain is climate, not economic.

A more recent group of works avoids the problem rather than solving it:
instead of first selecting the basis dimension and then, separately, the
smoothing parameter, they formulate both decisions as a single Bayesian
inference problem. \cite{sousa_bayesian_2024},
\cite{wakayama_horseshoe_2024}, \cite{dacruz_fast_2026}, and
\cite{tao_integrated_2025} use shrinkage priors (horseshoe,
spike-and-slab, sparse priors via variational inference) that select
which basis functions to retain and how much to penalize them within the
same estimation step. This unification solves the practical problem of
having to tune two parameters separately, but it does not answer PR3 in
the sense that matters to this thesis: by merging both decisions into a
single mechanism, it does not allow isolating how much of the result is
explained by the chosen basis family and how much by the smoothing
regime applied to it.

Of the 21 direct antecedents of PR3, only one works with economic or
financial series: \cite{rice_exploring_2023}, on the functional
volatility of crude-oil futures. But its form of regularization is not
the roughness penalization discussed by the other works in this section;
it is a non-negativity constraint on data-driven bases. No antecedent in
this corpus applies to economic or financial series the kind of explicit
separation between basis family and smoothing regime that does exist,
within a single family, in \cite{aguilera_comparative_2013} and
\cite{gao_factor_2024}. Section~\ref{sec:soa_economics} returns to this
point.

%---------------------------------------------------------
\section{Wavelet bases and non-smooth processes}
\label{sec:soa_wavelets}

PR4 asks how wavelet bases perform relative to polynomial or
trigonometric bases for non-smooth or non-stationary processes, but
among the 14 direct antecedents of this axis, applications that actually
use a wavelet basis are a minority: only \cite{amato_unsupervised_2025}
and \cite{yang_feature_2022} adopt one as the central component of their
method. Most of the works that address irregular processes respond to
that irregularity without resorting to wavelets at all, which is itself
relevant to PR4: within the recent practice documented by this corpus,
wavelets are not the default response to the lack of smoothness that
theory would suggest.

When the direct comparison is made, the result does not favor wavelets
systematically. \cite{amato_unsupervised_2025} find a considerable
advantage of wavelet-based clustering over B-spline bases (an adjusted
Rand index of 0.901 against 0.320), but \cite{salvatore_exploring_2016}
report the opposite: a wavelet-PCA analysis revealed no temporal changes
beyond those already captured by a Fourier-based FPCA with smoothing.
\cite{eslami_spline_2024} describe splines and wavelets as
``complementary,'' with similar mean-squared-error performance in one of
its case studies, and explicitly notes that the literature lacks
comprehensive comparisons between the two families. The pattern is the
same as in Section~\ref{sec:soa_basis_choice}: no single family
dominates universally, and the outcome depends on the type of
irregularity present in the data.

The most frequent strategy among the antecedents of PR4 is not to
replace the smooth basis with a wavelet, but to adapt the smooth basis
to the irregularity. \cite{jiao_enhanced_2022} align the location of the
basis functions with the change point they aim to detect, instead of
fixing their position in advance, and show that this alignment
substantially improves detection power relative to a standard FPCA.
\cite{horvath_time_2020} replace B-spline interpolation with cubic
Hermite polynomials for commodity forward curves, which exhibit breaks
at contract expiry points, and find a better fit with this alternative,
also non-wavelet. \cite{aguilera_comparative_2013} and
\cite{gao_factor_2024}, already discussed in
Section~\ref{sec:soa_smoothing}, absorb roughness through the
penalization parameter within the same B-spline family, rather than
switching families. None of these three strategies uses wavelets; all
three modify a smooth basis to accommodate irregularity.

None of the PR4 antecedents that directly compare wavelets against other
families works with economic or financial series. The two antecedents
of this axis that do -- \cite{aue_generalized_2017}, on intraday returns
of the SPY ETF, and \cite{rice_exploring_2023}, on the volatility of
crude-oil futures -- do not offer that comparison: the former mentions
wavelets only as a typical alternative, untested in the study itself,
and the latter regularizes through non-negativity constraints, not
through a wavelet basis. The question posed by PR4 therefore has no
direct empirical answer for the economic-financial domain within the
reviewed literature -- another point that Section~\ref{sec:soa_economics}
takes up.

%---------------------------------------------------------
\section{FDA in economics and finance; Latin American and Ecuadorian
antecedents}
\label{sec:soa_economics}

Functional data analysis is well established as an applied tool in
economics and finance: the 35 items in this axis cover electricity
demand, stock indices, exchange rates, commodities, cryptocurrencies,
and income distribution. But this breadth of domains contrasts with
methodological depth: only 11 of the 35 (31\%) reach the level of a
direct antecedent ($Q3\geq1$ and $Q4\geq1$), against 48\% of direct
antecedents in the full corpus. Most of the remaining 24 applications
use FDA as an analysis tool without discussing the choice of basis or
evaluating the quality of the resulting representation -- the same
pattern of incidental use already documented as a reason for exclusion
(E1) during Phase 5 screening -- with scores concentrated between 0 and
4 out of 8.

Among the 11 direct antecedents, none applies to economic or financial
series the separation between the effect of the basis family and the
effect of smoothing that does exist, within a single family, in
\cite{aguilera_comparative_2013} and \cite{gao_factor_2024}
(Section~\ref{sec:soa_smoothing}). \cite{shackleton_comparing_2024}
theoretically justify their choice of B-splines by the aperiodicity of
the realized volatility they study, but do not compare it against
another family. \cite{rodriguezcuadro_study_2025}, on the Colombian
stock market, combine P-splines with FPCA and clustering to
characterize the functional correlation between 26 companies and the
price of Brent crude oil, but likewise do not compare basis families
against each other. Each economic antecedent satisfies some subset of
the four valuation criteria -- theoretical justification, reported
parameters, evaluation of the representation -- but none satisfies all
four at once the way \cite{beyaztas_comparison_2022} or
\cite{aguilera_comparative_2013} do outside the economic domain.

Regional coverage is even more limited. Of the 91 items in the corpus,
only three have any connection to Latin America: \cite{marin_analyzing_2023},
on electricity demand in Colombia; \cite{rodriguezcuadro_study_2025},
already cited; and \cite{padilla2020} itself, on demographic and
economic indicators for the region. None of the 91 items retrieved by
the systematic search works specifically with Ecuadorian data. This
scarcity is not an artifact of the search strings -- Phase 4 had already
anticipated a low volume for the regional search string -- but a pattern
that is now confirmed within the set already screened and scored for
methodological quality: Latin American literature on FDA exists, but it
is scarce, and its intersection with the comparison of basis families
is, within this corpus, empty.

%---------------------------------------------------------
\section{Critical synthesis and the research gap}
\label{sec:soa_gap}

Before restating it as an actionable gap, the candidate gap must be
checked against the extraction matrix itself rather than asserted from
an impression of the five axes just reviewed. Cross-referencing the 25
items that explicitly compare two or more basis families against the
questions they address and their application domain narrows the picture
considerably: only 5 of those 25 touch PR5, applied to electricity
prices \cite{lisi_joint_2024}, wastewater epidemiology (a marginal
connection to PR5) \cite{salvatore_exploring_2016}, natural gas flows
\cite{chen_review_2021}, intraday returns of the SPY exchange-traded
fund \cite{aue_generalized_2017}, and commodity forward curves
\cite{horvath_time_2020}. None of these five isolates the effect of the
basis family from the effect of the smoothing or penalization regime as
an explicit methodological contribution; each reports model comparisons
or extensions rather than a design that controls both factors
separately. The only item in the corpus that does separate basis family
from estimation method explicitly, \cite{beyaztas_comparison_2022},
applies to weekly climate data, not to economic or financial series.
And, consistent with the regional scarcity documented in
Section~\ref{sec:soa_economics}, none of the 91 items applies to Latin
American or Ecuadorian data. The candidate gap therefore holds, but it
must be stated more precisely: the issue is not that systematic basis
comparison is absent from FDA in general -- 25 of 91 items (27.5\%)
perform one, and at least one isolates family from estimation method --
but that this comparison has not been made for volatile economic or
financial series, still less with the effect of smoothing isolated or
with Latin American data.

Three observations follow from the five axes reviewed above.
\textbf{What is established.} Representing functional data through
basis expansion is well established (B-splines dominate with 30 of 91
single-family uses, followed by empirical/FPCA bases with 12 of 91), and
the criteria for selecting basis dimension are mature and in widespread
use (GCV, AIC/BIC, and cross-validation appear in most items that report
a criterion). \textbf{What is fragmentary.} Systematic comparison across
basis families is not uncommon in the general FDA literature -- 25 of 91
items (27.5\%) explicitly compare two or more families, and at least one
(\cite{beyaztas_comparison_2022}) formally isolates the effect of the
basis family from the effect of the estimation method. But this
comparative literature concentrates on domains outside economics and
finance (climate, genetics, spectroscopy, geodesy, EEG): only 5 of those
25 works touch economic or financial series, and none of them separates
the effect of the family from the effect of the smoothing regime.
\textbf{What is absent.} No work in the reviewed literature offers a
systematic, controlled comparison of functional basis families for FPCA
representation of volatile economic or financial series that isolates
the effect of basis choice from the effect of penalization or the
smoothing regime, and no application of this kind exists for Latin
American or Ecuadorian data. This is precisely where \cite{padilla2020}
-- the article that motivates this thesis -- stops: it uses B-splines
exclusively, without comparison and without penalization, and its own
proposed future work explicitly calls for ``considering different
approximation basis, specifically for the case of GDP-GR'' -- textual
confirmation, from within the field's own literature, of the exact gap
this thesis addresses.

%---------------------------------------------------------
\section{Positioning of this thesis}
\label{sec:soa_positioning}

This thesis is designed to close the gap identified above on its
narrowest and most defensible terms: a controlled comparison of
functional basis families for volatile economic time series that
isolates the effect of basis choice from the effect of the smoothing
regime, applied to Latin American data. Chapter~\ref{chap:methods}
extends the dataset of \cite{padilla2020} -- the study that motivates
this work -- to an annual panel of GDP growth rates for 20 Latin
American countries from 1960 to 2024, incorporating six years absent
from the original study, among them the 2020 COVID-19 contraction and
the region's heterogeneous post-pandemic recovery. Unlike the 25
comparative works identified in Section~\ref{sec:soa_gap}, none of which
applies to economic or financial data with the effect of smoothing
controlled, the comparison here fixes reconstruction quality across
families before comparing them: the wavelet threshold is calibrated, per
country, to match the RMSE achieved by the B-spline fit, so that
differences in sparsity and in the resulting FPCA are attributable to
the basis family rather than to an uncontrolled difference in smoothing
level. This design answers, for the economic-financial domain that
Section~\ref{sec:soa_economics} found empty, the same question that
\cite{aguilera_comparative_2013} and \cite{gao_factor_2024} answer only
within a single basis family (Section~\ref{sec:soa_smoothing}): what is
the effect of the basis family once the smoothing regime is no longer
allowed to vary freely across it.

The scope of this comparison is deliberately narrower than some of the
alternatives reviewed in this chapter. It does not estimate a
data-driven orthonormal basis, as the matrix-completion and splinet
methods of Section~\ref{sec:soa_basis_choice} do, nor does it unify
basis selection and penalization into a single Bayesian mechanism, as
the shrinkage-prior methods of Section~\ref{sec:soa_smoothing} do; both
choices would reintroduce, in a different form, the confound this thesis
is designed to remove. It restricts the comparison to three classical,
non-data-driven families -- B-spline, Fourier, and wavelet (Haar, DB2,
Sym4) -- and evaluates them jointly on reconstruction fidelity,
coefficient economy, and the interpretability of the resulting
functional principal component analysis, directly extending the
single-basis analysis of \cite{padilla2020} in the three directions its
own proposed future work calls for. Within the literature reviewed in
this chapter, it is, to the best of this review's knowledge, the first
systematic comparison of this kind applied to Latin American economic
data.
